\chapter{Conclusion}\label{ch:conclusion}

This thesis has demonstrated that language design can meaningfully address the main challenges faced by smart contract developers and auditors today. Instead of placing additional burden on developers, we move responsibility to the compiler and type system. Vulnerabilities can be prevented by adopting a safety-by-default policy, while allowing developers to opt in to unsafe behavior on a case-by-case basis. Importantly, developers must be able to opt in to \textit{exactly only the behaviour they require}, nothing more. Further, increasing auditability of code is an important aspect to the language design in this ecosystem, due to the prevalence of third-party audits, especially for well-known projects. The fact that a large proportion of high-profile exploits involved audited code serves to demonstrate that there is significant potential for improvement in this area. Lastly, it is also important to provide developers with a sufficiently expressive set of features to natively model their desired behaviour, such that its type-safety is guaranteed at compile time, as opposed to manual re-implementations which are more prone to errors.

\section{Summary of Results}
As we have seen, the largest class of vulnerabilities that are amenable to language-level solutions are reentrancies. The majority of real-world examples would have been prevented simply by rejecting all reentrant calls by default -- this clearly shows a lack of understanding or tooling in the current ecosystem. Furthermore, it makes more sense for the compiler to preempt issues related to EVM intricacies -- otherwise developers are expected to know about and proactively prevent vulnerabilities that do not stem from their own code, but rather from EVM-specific behaviour. To solve these issues, we have introduced a system of path-specific reentrancy guards, which rejects all reentrant calls by default. When a developer requires a reentrant call, they can specifically opt in to exactly one reentrant call path. This is a much safer approach when compared to the standard practice of leaving an entrypoint entirely unguarded. As we have demonstrated, the overhead for these extra protections is relatively benign, with the single exception of read-only functions. However, the term ``read-only reentrancy'' has been coined specifically to describe the vulnerability this is preventing, demonstrating that this is a significant threat which is worth protecting against.
\par
Most importantly, we have clearly demonstrated that our approach to reentrancy protection can effectively prevent real-world vulnerabilities. Of the 54 examples of reentrancies in the dataset, 53 resulted from errors in developer code, and 16 of those projects did actually make use of reentrant calls in some way. Of those 16 vulnerabilities, all of them could have been prevented by selectively re-enabling only the required functionality. Of course, this does not rule out the potential for developer error -- it is still possible for a developer to explicitly allow a reentrant call path that is unsafe. However, given the vast reduction in the number of possibilities there are to reason about, it becomes much easier to find or rule out reentrancy vulnerabilities.
\par
Our second key feature is the introduction of side-effect annotations. It has two key benefits -- firstly, it enforces an ``effect-oriented'' approach to smart contract development, which is already relatively widespread among popular projects. This approach also serves as a mechanism which forces the developer to double-check their work, at least when it comes to critical state modifications. More importantly, the intent is to accelerate the auditing process, as well as enabling auditors to rule out certain interactions more quickly, allowing them to dive deeper into the space of possibilities without worrying about surface-level bugs. It also offloads the responsibility of correctly inferring which effects happen where to the compiler, rather than being prone to human errors. The developer overhead induced by these annotations is somewhat significant for very small contracts, but as contract size grows it becomes more negligible.
\par
Lastly, we have introduced the notion of algebraic data types to the EVM. It is a natural fit for multiple reasons - first of all, the modelling of state machines is a common requirement, and ADTs model these in a natural and type-safe way. Through exhaustive pattern matching, the compiler can also ensure that all possible cases are handled, eliminating the chance for developer error. ADTs such as Rust's \rustinline{Option} and \rustinline{Result} allow native handling of failures or errors, which is a major improvement over ``silent'' failure. Further, combining ADTs with the coding style of idiomatic Rust also naturally separates \textit{pure} functionality from that of the contract itself. In other EVM languages, functions are almost always defined in ``contract space'', meaning they cannot logically be separated from the underlying contract, even if the functionality is free of side effects. By contrast, we have found it more natural to define methods on the types themselves, meaning the code in ``contract space'' becomes sparser and more focused on the contract's intended functionality, as opposed to the implementation of arithmetic operations or other more abstract functionalities. This also paves the path for more natural code reuse, since the user-defined types can be fully separated from the smart contracts. Finally, this separation also enables the possibility of unit testing such pure code without the need of mocking a blockchain, contracts, or any such complexities.

\section{Limitations}
The increased complexity of the type system and compiler is concerning -- especially given the difficulty of type inference involving generics, method resolution, and further challenges. The two-level hierarchy of compilers further increases the risk of errors -- not only do we rely on our own compiler's correctness, but now also on that of the Solidity compiler. Of course, we only rely on relatively simple functionality, and the Solidity compiler is by far the most battle-tested option available. For the moment, the advantages provided by leveraging the tools available for Solidity code far outweigh the risk of compiler errors, but reducing the reliance on external tools does provide a natural way forward. \\
Regarding the complexity of the type system, it is clear that it does provide additional sources of compiler error. However, reasoning about such abstract types is a well-understood field of study, and Hindley-Milner type inference, as well as the specific approach we have based our inference algorithm on, are known to be sound. The intent is mostly to allow abstractions of functionality over monadic types, for example by defining addition over Rust's \rustinline{Option} type. Contract code cannot directly contain generics, only instantiate generic code with concrete types.

\section{Future Directions}

\subsection{Adoption Road-Blocks}
Before the language could be considered production-ready, several engineering milestones remain. Many intrinsic EVM features, such as event logging, delegate-calls, and programmatic contract creation are not yet supported. Further, the only tooling available is due to the reliance on the Solidity compiler, if this were eventually phased out, then tools such as debuggers, testing and deployment frameworks would need to be created. Additionally, the modularity of the language needs to be expanded to allow for libraries, contract composition, and other mechanisms for code reuse.

\subsection{Compiler Work}
A natural way forward is to reduce the reliance on the Solidity compiler. This can be done incrementally without any problem, since one could smoothly transition from generating native Solidity to including more and more inline assembly. Note that ``inline assembly'' in Solidity really just means inline Yul, which is the intermediate representation used by the compiler. Once the generated code consists only of this inline assembly, it can be easily transitioned to generate full Yul contracts instead. This would still benefit from the Solidity compiler's powerful optimizer stage, which operates on the Yul IR directly -- thus having comparable performance to Solidity.

\subsubsection{Interpreter}
Another interesting avenue would be the implementation of an interpreter for the language. This would provide an easy way of emulating the behaviour of internal functions, without the need for mocking the entire EVM or blockchain. Current debugging and testing tools have to deploy a local test-chain to run the contracts in, which is clearly a huge amount of overhead. While this isn't prohibitive for integration tests, unit testing is too slow for rapid prototyping work, and debugging smart contract code is quite tedious. Moreover, fuzz testing is heavily reliant on performance, and thus being able to directly interpret the code without the overhead of the EVM could provide a significant advantage over existing tools. If such an interpreter could also easily model external contract interactions, it may be possible to eliminate the mocking of the EVM entirely.

\subsection{User Studies}
One important open question is whether the effect system, as currently designed, strikes the right balance between safety and usability. While the annotations provide clear benefits for auditors, they also introduce significant developer overhead, particularly for smaller contracts. User studies would be a natural way to evaluate this tradeoff empirically. For instance, one could study whether auditors are able to find or rule out vulnerabilities more quickly in code written in our language with and without effect annotations, and compared to equivalent Solidity or Vyper code. Similarly, a study of developer experience could assess whether the annotation syntax is intuitive in practice, or whether alternative ways of declaring effects would reduce friction without sacrificing clarity. Such studies would also provide valuable data on the learning curve associated with adopting new languages in a high-stakes environment, and could directly inform future iterations of the language design. This kind of empirical evaluation is well-established in programming language research \cite{Stefik2013} \cite{Crichton2021}, especially when considering advanced type systems \cite{Ferdowsi2023}, and would complement the formal and dataset-based evaluation presented in this thesis.

\subsection{Language Extensions}
There are several interesting avenues for future language design research. For one, many systems involve multiple contracts interacting with one another. Because they all belong to the same project, the calls between the different contracts are trusted. As such, the effect system could be extended to model not only the current contract's storage, but also that of other known contracts. Furthermore, reentrancy guards could be omitted on one contract if it can only ever be accessed through another, leading to substantial gas savings. In a similar vein, many contracts end up being split into multiple parts, due to the size restriction imposed on contract bytecode. This could be done automatically by the compiler, perhaps even optimizing for the fewest number of cross-contract calls or some other cost metric, instead of relying on developers to figure out their own workarounds.
\par
Access control is another aspect of smart contracts that can be improved by language design -- many experimental languages introduce their own novel access control primitives. Even still, access control bugs result in a significant number of exploits, again demonstrating significant potential for improvement. However, it is difficult to define unifying principles that are shared between the large variety of access control mechanisms used by smart contracts in practice.
\par
In the last year or two, there has been a recent trend of exploits taking advantage of precision loss, even in audited projects. This suggests that existing tools provide insufficient support for reasoning about rounding errors. While it may be difficult for this to be addressed directly with the design of a language, perhaps other approaches such as static analysis could allow checking safe bounds on computations involving rounding errors.
\par
Lastly, once the type and effect systems have been stabilized, it becomes important to ensure their correctness. This could be done for instance by formally verifying their implementations, providing a foundation for a formally verified compiler stack.

\section{Closing Remarks}
The prevalence of reentrancy vulnerabilities is not an unavoidable consequence of the EVM, but a lack of adequate protection mechanisms in available tools. Furthermore, our language design in principle allows static analysis (or similar reasoning) to propagate across reentrant calls, which is not natively possible in any other existing languages. It also improves reasoning capabilities regarding side effects, which may help auditors to find or rule out vulnerabilities more easily. The coding style incentivized by the language also seems to naturally separate effectful code from more pure, abstract functionalities. We have also prioritized expressiveness in the type system to offload more responsibility to the compiler, so that developers and auditors can focus more on the intended functionality rather than worrying about compiler or EVM intricacies.
\par
The smart contract ecosystem has long treated security as an afterthought, relying on developer discipline and post-hoc auditing and analysis to catch vulnerabilities that should never have been possible in the first place. This thesis has argued, and empirically demonstrated, that language design offers a more principled path forward -- one where safety guarantees are enforced by construction rather than convention. While significant work remains before such a language could see widespread adoption, the convergence of independent efforts such as Fe's recent redesign toward similar principles suggests that the ecosystem is beginning to recognise this. The question is no longer whether language design can improve smart contract security, but how far that improvement can go.
